\documentclass[a4paper]{article}

\usepackage{graphicx}
\usepackage{amsmath,amssymb}
\usepackage{minted}
\usepackage{multirow}
\usepackage{float}
\usepackage{todonotes}
\usepackage{forest}
\usepackage{xstring}
\usepackage{xcolor}
\usepackage[most]{tcolorbox}
\usepackage{xparse}
\usepackage{natbib}

\usetikzlibrary{graphs,graphdrawing}
\usegdlibrary{layered}

% for external graphviz
\input{/home/donkarlo/Dropbox/repo/nd_ai_project/src/nd_ai/knoweldge_representation/language/formal/computer/programming/tex/latex/code_clips/external_graphviz.tex}

% for tags
\input{/home/donkarlo/Dropbox/repo/nd_ai_project/src/nd_ai/knoweldge_representation/language/formal/computer/programming/tex/latex/parts/control_sequence/kind/semtag/semtag.tex}

% for links
\input{/home/donkarlo/Dropbox/repo/nd_ai_project/src/nd_ai/knoweldge_representation/language/formal/computer/programming/tex/latex/code_clips/seehere.tex}

\title{Kalman Filter Equations}
\date{}

\begin{document}
    \maketitle
    
    The Kalman filter estimates the hidden state of a linear dynamical system from noisy measurements. It assumes that both the process noise and the measurement noise are Gaussian \cite{kalman-1960-a-new-approach-to-linear-filtering-and-prediction-problems}.
    
    
    \section{State-Space Model}
        
        The system is usually written as:
        
        \[
            x_k = A_k x_{k-1} + B_k u_k + w_k
        \]
        
        \[
            z_k = H_k x_k + v_k
        \]
        
        where:
        
        \[
            x_k
        \]
        
        is the hidden state of the system at time step \(k\).
        
        \[
            z_k
        \]
        
        is the measurement or observation at time step \(k\).
        
        \[
            u_k
        \]
        
        is the control input.
        
        \[
            A_k
        \]
        
        is the state transition matrix, which describes how the previous state moves to the current state.
        
        \[
            B_k
        \]
        
        is the control-input matrix.
        
        \[
            H_k
        \]
        
        is the observation matrix, which maps the hidden state into the measurement space.
        
        \[
            w_k
        \]
        
        is the process noise, usually assumed to be Gaussian:
        
        \[
            w_k \sim \mathcal{N}(0, Q_k)
        \]
        
        \[
            v_k
        \]
        
        is the measurement noise, also usually assumed to be Gaussian:
        
        \[
            v_k \sim \mathcal{N}(0, R_k)
        \]
        
        Here, \(Q_k\) is the process-noise covariance matrix, and \(R_k\) is the measurement-noise covariance matrix.
    
    
    \section{Prediction Step}
        
        In the prediction step, the filter predicts the next state before seeing the new measurement.
        
        The predicted state is:
        
        \[
            \hat{x}_{k|k-1} = A_k \hat{x}_{k-1|k-1} + B_k u_k
        \]
        
        The predicted covariance is:
        
        \[
            P_{k|k-1} = A_k P_{k-1|k-1} A_k^T + Q_k
        \]
        
        where:
        
        \[
            \hat{x}_{k|k-1}
        \]
        
        is the predicted state at time \(k\), based only on information available up to time \(k-1\).
        
        \[
            P_{k|k-1}
        \]
        
        is the predicted uncertainty of this state estimate.
        
        \[
            P_{k-1|k-1}
        \]
        
        is the previous corrected covariance.
    
    
    \section{Innovation}
        
        After the measurement \(z_k\) arrives, the filter computes the difference between the actual measurement and the predicted measurement.
        
        This difference is called the innovation or residual:
        
        \[
            y_k = z_k - H_k \hat{x}_{k|k-1}
        \]
        
        The innovation covariance is:
        
        \[
            S_k = H_k P_{k|k-1} H_k^T + R_k
        \]
        
        where \(y_k\) tells us how much the real measurement differs from what the model expected.
    
    
    \section{Kalman Gain}
        
        The Kalman gain decides how much the filter should trust the new measurement compared with the model prediction.
        
        \[
            K_k = P_{k|k-1} H_k^T S_k^{-1}
        \]
        
        or equivalently:
        
        \[
            K_k = P_{k|k-1} H_k^T
            \left(
                H_k P_{k|k-1} H_k^T + R_k
            \right)^{-1}
        \]
        
        If the measurement noise \(R_k\) is large, the Kalman gain becomes smaller, so the filter trusts the model more.
        
        If the predicted uncertainty \(P_{k|k-1}\) is large, the Kalman gain becomes larger, so the filter trusts the measurement more.
    
    
    \section{Correction Step}
        
        The corrected state estimate is:
        
        \[
            \hat{x}_{k|k} =
            \hat{x}_{k|k-1} + K_k y_k
        \]
        
        or:
        
        \[
            \hat{x}_{k|k} =
            \hat{x}_{k|k-1} +
            K_k
            \left(
                z_k - H_k \hat{x}_{k|k-1}
            \right)
        \]
        
        The corrected covariance is:
        
        \[
            P_{k|k} =
            \left(
                I - K_k H_k
            \right)
            P_{k|k-1}
        \]
        
        where:
        
        \[
            \hat{x}_{k|k}
        \]
        
        is the updated state estimate after using the measurement at time \(k\).
        
        \[
            P_{k|k}
        \]
        
        is the updated uncertainty after the correction step.
    
    
    \section{Compact Form}
        
        The Kalman filter can be summarized in two main phases.
        
        First, prediction:
        
        \[
            \hat{x}_{k|k-1} = A_k \hat{x}_{k-1|k-1} + B_k u_k
        \]
        
        \[
            P_{k|k-1} = A_k P_{k-1|k-1} A_k^T + Q_k
        \]
        
        Second, correction:
        
        \[
            K_k =
            P_{k|k-1} H_k^T
            \left(
                H_k P_{k|k-1} H_k^T + R_k
            \right)^{-1}
        \]
        
        \[
            \hat{x}_{k|k} =
            \hat{x}_{k|k-1}
            +
            K_k
            \left(
                z_k - H_k \hat{x}_{k|k-1}
            \right)
        \]
        
        \[
            P_{k|k} =
            \left(
                I - K_k H_k
            \right)
            P_{k|k-1}
        \]
    
    
    \section{Interpretation}
        
        The Kalman filter is a recursive estimator.
        It does not need the full history of measurements.
        At each time step, it only needs the previous estimate, the previous covariance, the current control input, and the current measurement.
        
        The prediction step answers:
        
        \[
            \text{What do I expect the state to be before seeing the new measurement?}
        \]
        
        The correction step answers:
        
        \[
            \text{How should I modify that expectation after seeing the new measurement?}
        \]
        
        The covariance matrix \(P\) represents uncertainty about the state estimate.
        The matrices \(Q\) and \(R\) control how uncertainty enters the system: \(Q\) describes uncertainty in the system dynamics, while \(R\) describes uncertainty in the sensor measurements.
        
        \bibliographystyle{plainnat}
        \bibliography{/home/donkarlo/Dropbox/repo/nd_shared_research/bibliography/bibliography.bib}
\end{document}